% =====================================================================
%  PROPOSAL P2 -- The security tax
%
%  Standalone. Set this as Overleaf's main document and compile with
%  pdfLaTeX; Overleaf runs BibTeX automatically. Locally:
%    pdflatex P2 && bibtex P2 && pdflatex P2 && pdflatex P2
% =====================================================================
\documentclass[11pt,a4paper]{article}
% =====================================================================
%  Shared preamble for the three research proposals.
%
%  Each proposal \input{}s this file and is otherwise standalone, so any
%  one of them can be set as Overleaf's main document and compiled on its
%  own. Nothing here is specific to a single proposal.
% =====================================================================
\usepackage[T1]{fontenc}
\usepackage[utf8]{inputenc}
% Times, matching the parent manuscript. mathptmx is in every TeX Live and on
% Overleaf; lmodern is deliberately not used, since it is absent from minimal
% distributions and this preamble should compile anywhere.
\usepackage{mathptmx}
\usepackage[a4paper,margin=25mm]{geometry}
\usepackage{amsmath,amssymb,amsthm}
\usepackage{booktabs}
\usepackage{tabularx}
\usepackage{xcolor}
\usepackage{titlesec}
\usepackage{enumitem}
\usepackage{microtype}
\usepackage[hidelinks,breaklinks]{hyperref}

% ---- shared style tokens, matching the parent manuscript ------------
\definecolor{netcol}{RGB}{31,78,121}
\definecolor{autocol}{RGB}{18,120,90}
\definecolor{bridgecol}{RGB}{176,96,0}
\definecolor{accent}{RGB}{140,20,20}

\newcommand{\layernet}{\textcolor{netcol}{\textbf{network}}}
\newcommand{\layerauto}{\textcolor{autocol}{\textbf{autonomy}}}
\newcommand{\dataset}[1]{\textsc{#1}}
\newcommand{\code}[1]{\texttt{#1}}
\newcommand{\Real}{\mathbb{R}}
\newcommand{\Norm}[1]{\left\lVert #1\right\rVert}
\newcommand{\MCR}{\mathrm{MCR}}

\newtheorem{claim}{Claim}
\newtheorem{proposition}{Proposition}
\newtheorem{definition}{Definition}
\theoremstyle{remark}
\newtheorem{remark}{Remark}

% ---- section styling -------------------------------------------------
\titleformat{\section}{\large\bfseries\color{netcol}}{\thesection}{0.7em}{}
\titleformat{\subsection}{\normalsize\bfseries}{\thesubsection}{0.6em}{}
\titlespacing*{\section}{0pt}{1.4em}{0.6em}

\setlist[itemize]{leftmargin=1.35em,topsep=0.3em,itemsep=0.25em}
\setlist[enumerate]{leftmargin=1.6em,topsep=0.3em,itemsep=0.25em}

% ---- a boxed "what this rests on" callout ---------------------------
\newenvironment{honestly}
  {\par\medskip\noindent\begin{tabular}{@{}p{2.5mm}@{\hspace{3mm}}p{0.9\linewidth}@{}}
   \textcolor{accent}{\rule{2.5mm}{\baselineskip}} &
   \small\itshape}
  {\end{tabular}\par\medskip}

% ---- the proposal header --------------------------------------------
\newcommand{\proposalheader}[3]{%
  \begin{center}
  {\Large\bfseries #1\par}\medskip
  {\large\itshape #2\par}\medskip
  {\small #3\par}
  \end{center}
  \medskip\hrule\medskip}


\begin{document}

\proposalheader
  {P2 --- The Security Tax}
  {Paying for Authenticity in Milliseconds: A Certified Budget for
   Post-Quantum Signing on the UAV Bus and Link}
  {Research proposal, RobustUAVs.ai project $\cdot$ Roger Nick Anaedevha $\cdot$
   drafted 2026-08-13\\
   Follow-on to \emph{From Wire to Flight} (NDSS submission)}

\begin{abstract}
\noindent
The composition theorem of the parent manuscript bounds the mission-completion
floor achievable when a residual budget $\Delta$ of undetected staleness reaches
the navigation input. The theorem never asks \emph{why} the input is stale. This
proposal takes that indifference seriously: post-quantum authentication on a
constrained UAV bus or command link imposes a latency that is, formally, the
same quantity an adversary spends. Security hardening and the attacker therefore
draw on one budget, and the certificate is what arbitrates between them. We
propose to measure the latency of classical and post-quantum schemes on
flight-controller-class hardware, model the fragmentation cost on DroneCAN and
MAVLink~2, and push the result through the existing chain to obtain a
\emph{certified authentication budget}: the maximum per-message signing latency
a platform can afford and remain airworthy.
\end{abstract}

% =====================================================================
\section{The gap}
% =====================================================================
Three facts collide, and nobody has put them in the same paper.

\paragraph{MAVLink~2 authenticates but does not encrypt.}
The protocol carries a 13-byte HMAC-SHA256 signature providing authentication
and integrity, and offers no payload confidentiality~\cite{trout_mavlink}. Its
adoption in practice is uneven and the community discussion of adding
encryption remains open.

\paragraph{DroneCAN has no authentication at all.}
The CAN protocol that the ArduPilot and PX4 projects use to reach peripherals
provides no origin authentication for bus
frames~\cite{dronecan_spec}. The parent project already ingests
$200{,}000$ real frames of flooding, fuzzing and replay against exactly this
surface~\cite{hcrl_uavcan}, so the attack side is characterised and the defence
side is empty.

\paragraph{Post-quantum migration is arriving on constrained platforms.}
Kyber-based group key agreement has been proposed for
swarms~\cite{pq_swarm_kyber}, NIST-standardised lattice
schemes~\cite{nist_pqc} are being integrated into MAVLink-based
systems, and post-quantum authentication for drone networks is under active
development~\cite{pq_drone_networks}; the broader UTM security picture is
surveyed in~\cite{utm_survey2026}.

Every one of those results reports the same metrics --- signature size, key
size, cycles, energy. \textbf{None reports what those milliseconds do to the
aircraft.} And the arithmetic is unforgiving on a bus whose payload is 8 bytes
in classical CAN and 64 in CAN-FD: a lattice signature of a few kilobytes
becomes multi-frame fragmentation, which is latency, which is staleness of
exactly the kind the parent manuscript's Lemma~1 already bounds.

% =====================================================================
\section{The claim}
% =====================================================================
\begin{quote}\itshape
Security hardening and the adversary draw on one budget, and the composition
theorem is what arbitrates between them.
\end{quote}

Formally, the residual budget $\Delta(\theta)\le H\min(\theta,s)$ of the parent
manuscript bounds \emph{adversarial} staleness. Authentication contributes
\emph{self-inflicted} staleness $\Delta_{\mathrm{auth}}$, a deterministic
function of the scheme, the message rate, and the transport MTU. The certificate
consumes their sum:

\begin{equation}
\label{eq:budget}
\rho(\theta) \;=\;
\big(\Delta(\theta) + \Delta_{\mathrm{auth}}\big)\,
\gamma\, e^{L T_c} \;\le\; m ,
\end{equation}

which, for a given corridor margin $m$, yields a \textbf{certified
authentication budget}: the largest $\Delta_{\mathrm{auth}}$ compatible with
airworthiness. Three corollaries follow immediately, and all are testable.

\begin{proposition}[Detector tightness buys cryptographic headroom]
\label{pr:trade}
Lowering $\theta$ shrinks $\Delta(\theta)$ and therefore enlarges the admissible
$\Delta_{\mathrm{auth}}$ at fixed $m$. Network-layer detection and cryptographic
hardening trade against one another through the corridor margin.
\end{proposition}

\begin{proposition}[A margin below which nothing certifies]
\label{pr:floor}
For each transport there exists a corridor margin below which no post-quantum
scheme in the candidate set satisfies Eq.~\eqref{eq:budget} at any detector
setting. The paper can name that margin per transport.
\end{proposition}

\begin{proposition}[Ranking by size is the wrong ranking]
\label{pr:ranking}
Because fragmentation is a step function of the transport MTU rather than linear
in signature bytes, the ordering of schemes under Eq.~\eqref{eq:budget} need not
agree with the ordering by signature size that the cryptographic literature
reports.
\end{proposition}

\begin{honestly}
Proposition~\ref{pr:trade} is the sentence the paper lives or dies on, and it is
only interesting if the numbers are real. A modelled $\Delta_{\mathrm{auth}}$
built from cited cycle counts would make the whole result a restatement of the
inputs. The measurement is the contribution; the theorem is the frame.
\end{honestly}

% =====================================================================
\section{What is reused, what is new}
% =====================================================================
\begin{center}
\small
\begin{tabularx}{\linewidth}{@{}XX@{}}
\toprule
\textbf{Reused unchanged} & \textbf{New} \\
\midrule
Theorem~1 and its proof ($\Delta$ is $\Delta$, whatever its source) &
$\Delta_{\mathrm{auth}}$: scheme $\times$ rate $\times$ MTU $\to$ seconds \\
Lemma~1, the residual-budget bound &
measured signing and verification latency on flight-controller-class hardware \\
\code{certificates/engine.py} and the corridor family &
the budget-allocation result and Propositions~\ref{pr:trade}--\ref{pr:ranking} \\
the \dataset{HCRL UAVCAN} ingest, 200k real frames~\cite{hcrl_uavcan} &
a defence side for the \code{c2\_link} and \code{intra\_bus} sublayers, which
the schema has but the corpus does not \\
\bottomrule
\end{tabularx}
\end{center}

% =====================================================================
\section{Experiment plan}
% =====================================================================
\subsection{Measure, do not cite}
Bench Ed25519, HMAC-SHA256 (the MAVLink~2 baseline), the NIST lattice signature
and KEM standards~\cite{nist_pqc}, and a compact-signature alternative, on two
targets: a flight-controller-class microcontroller (STM32H7, as on Pixhawk~6)
and a companion-computer-class board. Report sign latency, verify latency, and
--- the one that matters and the one nobody reports --- \emph{end-to-end
frame-to-frame} latency under a realistic message mix.

\subsection{Fragmentation model}
Map signature size to DroneCAN frame count to bus occupancy at realistic message
rates, yielding $\Delta_{\mathrm{auth}}$. Validate against a real bus, using the
timing of the benign portion of the \dataset{HCRL UAVCAN} capture as the
baseline.

\subsection{Compose}
Push $\Delta_{\mathrm{auth}}$ through the existing chain. Produce the certified
authentication budget as a function of corridor margin, and the scheme ranking
that Proposition~\ref{pr:ranking} predicts will differ from the size ordering.

\subsection{The adversarial interaction}
The most interesting case, and a genuinely new attack: an adversary who
\emph{induces} authentication work --- flooding the bus with frames carrying
invalid signatures --- turns the defence into the delay attack. Verification
cost is incurred before the signature is known to be bad, so the defender pays
per frame. This is directly measurable on the \dataset{HCRL} surface, and the
composition theorem bounds its effect on the mission.

\begin{honestly}
This attack is dual-use, and the write-up must handle it the way the parent
manuscript handled its evasion analysis: the mechanism is not novel to anyone
who has implemented signature verification, the result \emph{bounds} what the
attacker achieves, and the practical output is a defensive configuration
procedure. Rate-limiting and cheap pre-filters are the mitigation, and they
belong in the same section as the attack.
\end{honestly}

% =====================================================================
\section{Relation to the parent manuscript's limitations}
% =====================================================================
\begin{center}
\small
\begin{tabularx}{\linewidth}{@{}p{0.42\linewidth}p{0.10\linewidth}X@{}}
\toprule
\textbf{Limitation} & \textbf{Closed?} & \textbf{How} \\
\midrule
Analytical anchor and detector scope (one detector, a scalar operating point) &
\textbf{weakly} &
exercises a second attack surface, the bus, but still with a scalar operating
point \\
Evidence-class imbalance & no & unchanged \\
\bottomrule
\end{tabularx}
\end{center}

\noindent
P2's value is not limitation-closing; it is that it adds a defence side to two
sublayers the schema already models and the corpus already covers, and it does
so with data that needs no dataset gatekeeper.

% =====================================================================
\section{Platform integration}
% =====================================================================
A \textbf{Crypto Budget} page: choose a transport (MAVLink~2, DroneCAN, CAN-FD),
a scheme, a message rate, and a corridor margin; receive the certified floor and
a pass/fail against the $\MCR\ge0.90$ reference consistent with
DO-326A~\cite{do326a}. This is the most operationally useful page the platform
could carry, and it requires no dataset at all.

% =====================================================================
\section{Risks}
% =====================================================================
\begin{itemize}
\item \textbf{Hardware access.} A microcontroller development board is
inexpensive; a full flight controller is needed only for the bus validation, not
for the timing measurement. If neither is reachable, benchmarks on a pinned ARM
target are a defensible second best --- but the paper must say so, and the
provenance tag must say so too.
\item \textbf{Reviewers may read it as post-quantum benchmarking.} The framing
must lead with the budget result, not the numbers. The numbers are the
instrument, not the finding.
\item \textbf{Scheme selection may date quickly.} Pin the candidate set to the
standardised algorithms~\cite{nist_pqc} and report the model, so a new scheme
can be placed in the ranking without re-running the campaign.
\end{itemize}

% =====================================================================
\section{Venue and timeline}
% =====================================================================
Target ACM AsiaCCS, ESORICS, or IEEE EuroS\&P; NDSS if the induced-authentication
attack lands cleanly. Four to five months, and the bench work parallelises with
P1's measurement campaign.

\small
\bibliographystyle{plain}
\bibliography{proposals}

\end{document}
